% =====================================================================
\chapter{Coste, energía y límites físicos}
\label{tema:6}
% =====================================================================

Los temas anteriores han tratado el rendimiento de un dispositivo y de un conjunto de ellos. Este tema convierte esas magnitudes en las que deciden un proyecto real: duración, coste económico, energía y refrigeración. Después compara la energía de las operaciones con la del movimiento de datos y con el límite físico de Landauer para el borrado de información. Se siguen el §2.1 de Kaplan \emph{et al.}~\cite{kaplan}, el §4.1 y el apéndice~B de Chowdhery \emph{et al.}~\cite{palm}, la tarea~1 de CS336~\cite[§4.3]{cs336a1}, el trabajo de Horowitz~\cite{horowitz}, el capítulo~6 de Hennessy y Patterson~\cite{hp} sobre centros de datos y el artículo de Landauer~\cite{landauer}.

\section{Duración y coste de un entrenamiento}

\begin{proposicion}[Duración de un entrenamiento]\label{prop:duracion}
Sea un entrenamiento de $C$ FLOP en $n_{\mathrm{d}}$ dispositivos de rendimiento máximo $\Ppico$ con utilización $\MFU$. Su duración es
\[
t=\frac{C}{n_{\mathrm{d}}\,\Ppico\,\MFU},
\]
y consume $n_{\mathrm{d}}t$ dispositivos-segundo, que no dependen de $n_{\mathrm{d}}$ si el MFU se mantiene.
\end{proposicion}

El coste económico es, en primera aproximación, el número de GPU-hora por el precio de la GPU-hora. Duplicar el número de dispositivos reduce la duración a la mitad sin cambiar el coste, mientras el MFU no caiga; por la ley de Amdahl (\cref{teo:amdahl}) y el coste creciente de la comunicación, en la práctica cae.

\addtocounter{definicion}{1}%

\section{Energía, potencia y refrigeración}

\begin{definicion}[PUE]
La \emph{eficiencia en el uso de la energía} (PUE, \emph{Power Usage Effectiveness}) de un centro de datos es el cociente entre la energía total que consume y la que consumen los equipos de cómputo. Siempre es $\PUE\ge 1$; el exceso corresponde a refrigeración, conversión eléctrica y otros servicios~\cite[cap.~6]{hp}.
\end{definicion}

La energía de un entrenamiento es, aproximadamente, $E\approx(\text{GPU-hora})\cdot P_{\mathrm{GPU}}\cdot\PUE$, donde $P_{\mathrm{GPU}}$ es la potencia media por GPU (la máxima de la A100 SXM es de 400~W~\cite{a100}). Esta estimación ignora las CPU, la red y el almacenamiento.

Casi toda la energía eléctrica se disipa como calor, que debe evacuarse. Si un fluido de calor específico $c$ entra a una temperatura y sale $\Delta T$ más caliente, el calor que extrae por unidad de tiempo es
\[
\dot{Q}=\dot{m}\,c\,\Delta T,
\]
con $\dot{m}$ el caudal másico. Para el agua, $c\approx 4186$~J/(kg$\cdot$K); para el aire, $c\approx 1005$~J/(kg$\cdot$K), con una densidad unas 800 veces menor. Evacuar los $1024\cdot 400=409{,}6$~kW de las GPU del \cref{ej:chinchilla} con agua que se calienta 10~K exige $\dot{m}=409\,600/(4186\cdot 10)\approx 9{,}8$~kg/s, casi diez litros por segundo. Por eso los sistemas actuales de alta densidad pasan de la refrigeración por aire a la refrigeración líquida directa sobre el chip (\cref{eje:6-refrigeracion}).

\section{Operar frente a mover datos}

Horowitz~\cite{horowitz} recopiló la energía de operaciones elementales en una tecnología de 45~nm. Las cifras, de orden de magnitud, muestran que una suma en fp32 cuesta unos $0{,}9$~pJ y una multiplicación en fp32, unos $3{,}7$~pJ, mientras que leer 32~bits de DRAM cuesta unos $640$~pJ: mover el dato cuesta unas $640/3{,}7\approx 170$ veces más que operar con él. Las técnicas de los temas~1 a~3 (matrices sistólicas, \emph{tiling}, fusión, FlashAttention) ahorran energía por la misma razón por la que ahorran tiempo: reducen el tráfico con la memoria lejana.

\addtocounter{definicion}{1}%

\section{El límite de Landauer}

\begin{teorema}[Principio de Landauer]\label{teo:landauer}
Borrar un bit de información, es decir, llevar un sistema físico que almacena un bit desconocido (0 o 1 con igual probabilidad) a un estado fijo, en contacto con un entorno a temperatura $T$, disipa en ese entorno un calor
\[
Q\ \ge\ \kB T\ln 2,
\]
donde $\kB=1{,}380649\cdot 10^{-23}$~J/K es la constante de Boltzmann.
\end{teorema}

\addtocounter{definicion}{1}%

\begin{observacion}
El principio solo limita las operaciones \emph{irreversibles}, que destruyen información. Una suma $c=a+b$ que sobrescribe $a$ y $b$ es irreversible; la computación reversible, que conserva la información necesaria para deshacer cada paso, podría en principio operar por debajo de $\kB T\ln 2$ por operación. El vínculo con el módulo~6 es directo: $\kB\ln 2$ es la entropía de un bit expresada en julios por kelvin.
\end{observacion}

